\chapter{Fungsi Dua Variabel}\label{ch:b1:multivar}

Tahun ini berakhir dengan langkah pertama ke dimensi yang lebih tinggi: yaitu fungsi
$f(x, y)$ atas dua variabel real. Segalanya terampat --- limit,
\omterm{def:b1:continuity:continuous}{kekontinuan}, \omterm{def:b1:derivative:def}{turunan}, ekstremum --- tetapi setiap gagasannya memperoleh puntiran:
karena limitnya dapat didekati sepanjang setiap arah sekaligus, \omterm{def:b1:derivative:def}{turunannya}
terbelah menjadi \omterm{def:b1:multivar:partial}{turunan parsial}, dan \omterm{def:b1:multivar:partial}{gradiennya} menunjuk jalan mendaki. Adapun
teori lengkapnya (diferensial, $\R^n$ yang umum, submanifold) termasuk
tahun kedua; sedangkan di sini kita menegakkan kosakata dan teorema
jujur yang pertama.

\section{Bidang $\R^2$; kekontinuan}

\begin{definition}\label{def:b1:multivar:topology}
Pada $\R^2$, pakailah norma Euklides $\norm{(x,y)} = \sqrt{x^2 + y^2}$
(\cref{ch:b1:euclid}). Bola buka, \omterm{def:b1:topology:open}{persekitaran}, dan \emph{\omterm{def:b1:logic:sets}{himpunan} bagian
\omterm{def:b1:topology:open}{terbuka}} pada $\R^2$ didefinisikan persis seperti pada
\cref{ch:b1:topology}, dengan bola menggantikan \omterm{prop:b1:reals:intervals}{selang}. Sebuah fungsi
$f \colon U \to \R$ ($U \subseteq \R^2$ \omterm{def:b1:topology:open}{terbuka}) disebut \emph{\omterm{def:b1:continuity:continuous}{kontinu}
di} $a \in U$ ketika
\[
\forall\varepsilon > 0,\ \exists\delta > 0, \quad
\norm{X - a} \leq \delta \implies \abs{f(X) - f(a)} \leq
\varepsilon,
\]
dengan pencirian barisan yang sama seperti pada satu variabel.
Adapun jumlah, hasil kali, hasil bagi dan komposisi dengan fungsi
satu variabel yang \omterm{def:b1:continuity:continuous}{kontinu} mengawetkan \omterm{def:b1:continuity:continuous}{kekontinuan}; sedangkan \omterm{def:b1:logic:map}{pemetaan} \omterm{prop:b1:vspaces:coordinates}{koordinatnya}
\omterm{def:b1:continuity:continuous}{kontinu}, jadi demikian pula \omterm{def:b1:poly:def}{polinomial} pada $(x,y)$.
\end{definition}

\begin{example}[Batas kutub, cara yang bersih untuk membuktikan sebuah limit]
\label{ex:b1:multivar:polarbound}
Tunjukkan bahwa $f(x, y) = \dfrac{x^2y^2}{x^2 + y^2}$ (dengan $f(0,0) =
0$) \omterm{def:b1:continuity:continuous}{kontinu} di titik asalnya. Dalam \omterm{prop:b1:vspaces:coordinates}{koordinat} kutub $x =
\rho\cos\theta$, $y = \rho\sin\theta$:
\[
\abs{f} = \frac{\rho^4\cos^2\theta\sin^2\theta}{\rho^2}
= \rho^2\,(\cos\theta\sin\theta)^2 \leq \frac{\rho^2}{4}
\xrightarrow[\rho \to 0]{} 0 ,
\]
yaitu batas yang \emph{tak bergantung pada $\theta$}: sehingga apa pun arah
pendekatannya, nilainya terperas ke $0$. Keseragaman pada
$\theta$ itulah inti seluruhnya --- karena batas seperti $\abs g =
\abs{\cos\theta\sin\theta}$ (tanpa $\rho$ yang tersisa) tak membuktikan apa pun,
dan memang $g$ itulah perangkap radial pada
contoh berikutnya.
\end{example}

\begin{example}[Perangkap radialnya]\label{ex:b1:multivar:trap}
Misalkan $f(x, y) = \dfrac{xy}{x^2 + y^2}$ untuk $(x,y) \neq (0,0)$, dengan $f(0,
0) = 0$. Sepanjang setiap sumbunya, $f = 0 \to 0$; tetapi sepanjang diagonal $y
= x$, $f(x, x) = \frac12 \not\to 0$. \emph{Jadi tak ada limit di titik asalnya}:
karena mendekat sepanjang setiap garis, dan bahkan menemukan limit yang sama sepanjang
masing-masingnya, belumlah cukup (di sini limit garisnya tak sepakat; sedangkan contoh yang lebih buruk
sepakat sepanjang semua garis namun gagal sepanjang sebuah parabola,
\cref{exo:b1:multivar:3}). \omterm{def:b1:continuity:continuous}{Kekontinuan} pada setiap variabel secara terpisah
tak mengakibatkan \omterm{def:b1:continuity:continuous}{kekontinuan}.
\end{example}

\section{Turunan parsial}

\begin{definition}\label{def:b1:multivar:partial}
Adapun \emph{turunan parsial}\index{turunan parsial} $f$ di
$(a, b)$ adalah \omterm{def:b1:derivative:def}{turunan} satu variabel sepanjang sumbunya:
\[
\frac{\partial f}{\partial x}(a,b)
= \lim_{h \to 0} \frac{f(a + h, b) - f(a,b)}{h},
\qquad
\frac{\partial f}{\partial y}(a,b)
= \lim_{k \to 0} \frac{f(a, b + k) - f(a,b)}{k}.
\]
$f$ disebut berkelas $C^1$ pada $U$ ketika keduanya ada dan \omterm{def:b1:continuity:continuous}{kontinu} pada
$U$. Adapun \emph{gradien}\index{gradien} adalah $\nabla f(a,b) =
\bigl(\frac{\partial f}{\partial x},\, \frac{\partial f}{\partial
y}\bigr)(a,b)$.
\end{definition}

\begin{theorem}[$C^1$ mengakibatkan bidang singgung]\label{thm:b1:multivar:c1}
Misalkan $f$ berkelas $C^1$ pada $U$ dan $(a,b) \in U$. Maka, ketika $(h, k) \to
(0,0)$:
\[
f(a + h, b + k) = f(a, b)
+ h\,\frac{\partial f}{\partial x}(a,b)
+ k\,\frac{\partial f}{\partial y}(a,b)
+ o\bigl(\norm{(h,k)}\bigr) .
\]
Khususnya $f$ \omterm{def:b1:continuity:continuous}{kontinu}, dan grafik $z = f(x,y)$ mempunyai pada
setiap titiknya bidang singgung yang terbaca dari rumusnya.
\end{theorem}

\begin{example}[Hampiran linear dalam kerja]
\label{ex:b1:multivar:linapprox}
Taksirlah $f(1.02,\ 0.99)$ bagi $f(x, y) = x^3y^2$. Di $(1, 1)$:
$f = 1$, $\frac{\partial f}{\partial x} = 3x^2y^2 = 3$,
$\frac{\partial f}{\partial y} = 2x^3y = 2$, sehingga
\cref{thm:b1:multivar:c1} memberikan
\[
f(1.02,\ 0.99) \approx 1 + 3\,(0.02) + 2\,(-0.01) = 1.04 ,
\]
lawan nilai sejatinya $1.02^3 \times 0.99^2 = 1.04006\dots$ ---
sehingga galatnya berorde dua pada tambahannya, sebagaimana dijanjikan
$o(\norm{(h,k)})$. Adapun bidang singgung grafiknya di
$(1, 1, 1)$ adalah $z = 1 + 3(x - 1) + 2(y - 1)$, yaitu persamaan
yang tersirat pada taksirannya.
\end{example}

\begin{proof}
Geserlah satu \omterm{prop:b1:vspaces:coordinates}{koordinat} setiap kali:
\[
f(a+h, b+k) - f(a,b)
= \bigl[f(a+h, b+k) - f(a, b+k)\bigr] + \bigl[f(a, b+k) -
f(a,b)\bigr].
\]
Menurut teorema nilai rata-rata satu variabel
(\cref{thm:b1:derivative:mvt}), kurung pertamanya adalah $h\,
\frac{\partial f}{\partial x}(a + \theta h,\, b + k)$ bagi suatu
$\theta \in \intoo{0}{1}$, sedangkan yang kedua adalah $k\,\frac{\partial
f}{\partial y}(a,\, b + \theta' k)$. Lalu \omterm{def:b1:continuity:continuous}{kekontinuan} \omterm{def:b1:multivar:partial}{turunan parsialnya} di
$(a,b)$ memungkinkan kita menulis masing-masingnya sebagai (nilai di $(a,b)$) $+$
(galat $\to 0$); sehingga galat totalnya adalah $h\,\varepsilon_1 +
k\,\varepsilon_2 = o(\norm{(h,k)})$ karena $\abs h, \abs k \leq
\norm{(h,k)}$.
\end{proof}

\begin{theorem}[Kaidah rantai]\label{thm:b1:multivar:chain}
Misalkan $f$ berkelas $C^1$ pada $U$ dan $t \mapsto (x(t), y(t))$ berkelas $C^1$ dari
sebuah \omterm{prop:b1:reals:intervals}{selang} ke dalam $U$. Maka $g(t) = f\bigl(x(t), y(t)\bigr)$ berkelas
$C^1$, dengan\index{kaidah rantai!dua variabel}
\[
g'(t) = x'(t)\,\frac{\partial f}{\partial x}\bigl(x(t),y(t)\bigr)
+ y'(t)\,\frac{\partial f}{\partial y}\bigl(x(t),y(t)\bigr)
= \bigl\langle \nabla f,\ (x', y')\bigr\rangle .
\]
\end{theorem}

\begin{proof}
Terapkanlah \cref{thm:b1:multivar:c1} di $(x(t), y(t))$ dengan $(h, k) =
(x(t+s) - x(t),\, y(t+s) - y(t))$: karena ketika $s \to 0$, sifat \omterm{def:b1:derivative:def}{dapat diturunkan}
satu variabelnya memberikan $h = s\,x'(t) + o(s)$ dan $k = s\,y'(t)
+ o(s)$, sehingga $\norm{(h, k)} = O(s)$ dan
\[
g(t+s) - g(t)
= h\,\frac{\partial f}{\partial x} + k\,\frac{\partial
f}{\partial y} + o\bigl(\norm{(h,k)}\bigr)
= s\,\Bigl(x'\,\frac{\partial f}{\partial x} +
y'\,\frac{\partial f}{\partial y}\Bigr) + o(s) ,
\]
dengan galat terakhirnya menyerap baik kedua $o(s)$ milik $h$ dan $k$
(yang dikalikan nilai tetap \omterm{def:b1:multivar:partial}{turunan parsialnya}) maupun
$o(O(s))$ pada taksiran bidang singgungnya. Bagilah dengan $s$ lalu biarkanlah
$s \to 0$. Adapun \omterm{def:b1:continuity:continuous}{kekontinuan} $g'$ menyusul dari \omterm{def:b1:continuity:continuous}{kekontinuan} semua
bahannya.
\end{proof}

\begin{example}[Kaidah rantainya, diperiksa dengan dua cara]
\label{ex:b1:multivar:chainrun}
Misalkan $f(x, y) = x^2 y$ dan $g(t) = f(t, t^2)$. Secara langsung: $g(t) =
t^2\cdot t^2 = t^4$, sehingga $g'(t) = 4t^3$. Lewat kaidah rantainya:
$\frac{\partial f}{\partial x} = 2xy$ dan $\frac{\partial
f}{\partial y} = x^2$, yang dinilai sepanjang kurva $(t, t^2)$:
\[
g'(t) = 1\cdot 2t\cdot t^2 + 2t\cdot t^2 = 2t^3 + 2t^3 = 4t^3 .
\]
Kedua perhitungannya sepakat, dan pembelahannya bermakna: karena $2t^3$
pertumbuhannya datang dari bergerak \emph{ke kanan} lewat
kemiringan $x$, dan $2t^3$ dari bergerak \emph{ke atas} lewat
kemiringan $y$. Pada kurva yang tak berbentuk \omterm{def:b1:topology:closed}{tertutup} bagi $g$, hanya
perhitungan keduanya yang bertahan --- dan itulah inti teoremanya.
\end{example}

\begin{remark}[Membaca gradiennya]
Sepanjang arah satuan $u$, kaidah rantai pada $t \mapsto f(a + tu)$
memberikan \emph{\omterm{def:b1:derivative:def}{turunan} berarah} $\langle \nabla f(a),
u\rangle$: yang maksimal ketika $u$ menunjuk sepanjang $\nabla f(a)$ (menurut
Cauchy--Schwarz, \cref{thm:b1:euclid:cs}). Jadi \omterm{def:b1:multivar:partial}{gradiennya} adalah
arah pendakian tercuram, dan ia \omterm{def:b1:euclid:orthogonal}{ortogonal} terhadap kurva
aras $\{f = c\}$ (turunkanlah $f$ sepanjang sebuah kurva yang digambar di dalam
\omterm{def:b1:logic:sets}{himpunan} aras: maka kaidah rantainya memberikan $\langle\nabla f,\,
\text{singgung}\rangle = 0$).
\end{remark}

\begin{example}[Kurva aras dan gradien, pada satu fungsi]
\label{ex:b1:multivar:levelcurves}
Ambillah $f(x, y) = x^2 - y^2$. \omterm{def:b1:logic:sets}{Himpunan} arasnya: $\{f = c\}$ berupa
hiperbola yang \omterm{def:b1:topology:open}{terbuka} kiri-kanan untuk $c > 0$, atas-bawah untuk $c < 0$,
dan sepasang garis bersilangan $y = \pm x$ untuk $c = 0$ --- yaitu
peta kontur sebuah pelana gunung, dengan \omterm{met:b1:multivar:monge}{titik pelananya} di
titik asalnya, yaitu tempat kedua garis yang beraras nol bersilangan. \omterm{def:b1:multivar:partial}{Gradiennya}: $\nabla
f = (2x, -2y)$. Pada titik $(2, 1)$ (di aras $c = 3$):
$\nabla f = (4, -2)$, sedangkan vektor singgung kurva
arasnya, yang diparameterkan dekat titik itu oleh $\bigl(t,
\sqrt{t^2 - 3}\bigr)$, adalah $\bigl(1, \frac{t}{\sqrt{t^2 -
3}}\bigr) = (1, 2)$ di $t = 2$ --- dan memang
\[
\langle (4, -2),\ (1, 2)\rangle = 4 - 4 = 0 :
\]
yaitu \omterm{def:b1:multivar:partial}{gradien} yang tegak lurus konturnya, yang menunjuk ke nilai $f$
yang lebih tinggi (di sini: menjauhi sumbu $y$). Dua bacaan
lagi: bahwa \omterm{def:b1:multivar:partial}{gradiennya} lenyap tepat pada pelananya, tempat
peta konturnya menjepit; dan bahwa \emph{garis} singgung pada
kurva aras di $(2,1)$ adalah $4(x - 2) - 2(y - 1) = 0$, yaitu $2x -
y = 3$ --- yakni persamaan ``$\langle \nabla f,\ M - M_0\rangle =
0$'' yang mengampat garis singgung elips pada
\cref{exo:b1:curves:11}.
\end{example}

\begin{theorem}[Schwarz]\label{thm:b1:multivar:schwarz}
Jika $f$ berkelas $C^2$ (yakni \omterm{def:b1:multivar:partial}{turunan parsial} atas \omterm{def:b1:multivar:partial}{turunan parsialnya} ada dan
\omterm{def:b1:continuity:continuous}{kontinu}), maka\index{teorema Schwarz}
\[
\frac{\partial^2 f}{\partial x\,\partial y}
= \frac{\partial^2 f}{\partial y\,\partial x} .
\]
\end{theorem}
\admitted

\section{Ekstremum lokal}

\begin{method}[Kajian ekstremum, tertata]
\label{met:b1:multivar:extremamethod}
\begin{enumerate}
  \item \emph{Pecahkanlah $\nabla f = 0$ selengkapnya.} Faktorkanlah setiap
        \omterm{def:b1:multivar:partial}{turunan parsialnya} kapan pun mungkin (karena hasil kali faktor \omterm{def:b1:linmaps:def}{linear}
        membelah sistemnya menjadi kasus yang bening, seperti pada
        \cref{ex:b1:multivar:fourpoints} di bawah); dan kasus yang terlupa
        adalah \omterm{thm:b1:multivar:critical}{titik kritis} yang terlupa.
  \item \emph{Golongkanlah setiap titiknya} dengan data Monge $r, s,
        t$ --- yang dihitung ulang \emph{pada setiap titiknya}, tak pernah sekali untuk
        semuanya.
  \item \emph{Jika $rt - s^2 = 0$}, periksalah $f$ secara langsung sepanjang
        kurva yang terpilih baik lewat titiknya (garis dahulu, lalu
        parabola), sembari memburu entah dua tanda (jadi tak ada ekstremum)
        atau satu tanda yang terkunci beserta argumen yang meliputi semua
        arahnya.
  \item \emph{Mundurlah untuk gambaran globalnya}: periksalah
        perilakunya di takhingga (karena minimum lokal boleh jadi bukan yang global),
        dan jika daerah asalnya tak \omterm{def:b1:topology:open}{terbuka}, perlakukanlah batasnya
        secara terpisah (\cref{exo:b1:multivar:12}) --- karena teorema titik
        kritisnya hanya melihat titik dalamnya.
\end{enumerate}
\end{method}

\begin{theorem}[Titik kritis]\label{thm:b1:multivar:critical}
Jika $f$ ($C^1$ pada \omterm{def:b1:topology:open}{himpunan terbuka} $U$) berekstremum lokal di $(a,b)
\in U$, maka $\nabla f(a,b) = (0,0)$: jadi titiknya bersifat
\emph{kritis}\index{titik kritis}.
\end{theorem}

\begin{proof}
Fungsi satu variabel $x \mapsto f(x, b)$ dan $y \mapsto f(a,
y)$ berekstremum lokal di dalam, di $a$ dan $b$ berturut-turut:
sehingga \cref{prop:b1:derivative:fermat} membunuh kedua \omterm{def:b1:multivar:partial}{turunan parsialnya}.
\end{proof}

\begin{method}[Uji orde dua (notasi Monge)]\label{met:b1:multivar:monge}
Pada sebuah \omterm{thm:b1:multivar:critical}{titik kritis} fungsi $C^2$, tetapkanlah
\[
r = \frac{\partial^2 f}{\partial x^2},
\qquad
s = \frac{\partial^2 f}{\partial x \partial y},
\qquad
t = \frac{\partial^2 f}{\partial y^2}
\qquad (\text{nilai pada titiknya}).
\]
\begin{itemize}
  \item Jika $rt - s^2 > 0$: maka ada ekstremum lokal --- yakni minimum untuk $r > 0$,
        dan maksimum untuk $r < 0$;
  \item jika $rt - s^2 < 0$: maka tak ada ekstremum (melainkan sebuah \emph{titik
        pelana}\index{titik pelana});
  \item jika $rt - s^2 = 0$: maka ujinya bungkam; periksalah secara langsung.
\end{itemize}
(Pembenarannya --- yaitu penjabaran Taylor--Young pada orde $2$ dan
kajian tanda bentuk kuadratik $r h^2 + 2shk + tk^2$ ---
dijalankan pada tahun kedua; sedangkan di sini ujinya dipakai sebagai perkakas
kerja.)
\end{method}

\begin{example}\label{ex:b1:multivar:extrema}
$f(x,y) = x^3 + y^3 - 3xy$. Titik kritisnya: $\nabla f = (3x^2 -
3y,\; 3y^2 - 3x) = 0$ memberikan $y = x^2$ dan $x = y^2$, sehingga $x = x^4$:
jadi $x \in \{0, 1\}$: yaitu titik $(0,0)$ dan $(1,1)$.

\omterm{def:b1:derivative:def}{Turunan} keduanya: $r = 6x$, $s = -3$, $t = 6y$. Di $(0,0)$: $rt -
s^2 = -9 < 0$: jadi pelana. Di $(1,1)$: $rt - s^2 = 36 - 9 > 0$, $r = 6
> 0$: jadi minimum lokal, $f(1,1) = -1$. (Yang tak global: karena $f(x, 0) = x^3 \to
-\infty$.)
\end{example}

\begin{example}[Kajian empat titik, selengkapnya]
\label{ex:b1:multivar:fourpoints}
$f(x, y) = xy\,(3 - x - y) = 3xy - x^2y - xy^2$. \omterm{def:b1:multivar:partial}{Gradiennya}:
\[
\frac{\partial f}{\partial x} = y\,(3 - 2x - y),
\qquad
\frac{\partial f}{\partial y} = x\,(3 - x - 2y).
\]
Titik kritisnya: jika $y = 0$, maka persamaan keduanya memberikan $x \in
\{0, 3\}$; jika $x = 0$, maka yang pertama memberikan $y \in \{0, 3\}$; sedangkan jika $xy
\neq 0$, pecahkanlah $2x + y = 3$, $x + 2y = 3$: jadi $x = y = 1$. Empat
titik: $(0,0)$, $(3,0)$, $(0,3)$, $(1,1)$. \omterm{def:b1:derivative:def}{Turunan} keduanya:
$r = -2y$, $s = 3 - 2x - 2y$, $t = -2x$.
\begin{itemize}
  \item $(1,1)$: $r = -2$, $s = -1$, $t = -2$: $rt - s^2 = 3 >
        0$, $r < 0$: jadi maksimum lokal, $f(1,1) = 1$.
  \item $(0,0)$: $r = t = 0$, $s = 3$: $rt - s^2 = -9 < 0$:
        jadi pelana; demikian pula $(3, 0)$ ($s = -3$) dan $(0, 3)$: jadi tiga
        pelana.
\end{itemize}
Maksimumnya hanya lokal: $f(-T, -T) = T^2(3 + 2T) \to
+\infty$. Periksa kesetangkupannya: $f(x, y) = f(y, x)$, dan memang
\omterm{def:b1:logic:sets}{himpunan} kritis dan penggolongannya setangkup pada $x
\leftrightarrow y$. Tafsirannya: di antara persegi \omterm{def:b1:curves:arclength}{panjang} berkelonggaran
$x, y \geq 0$, $x + y \leq 3$, hasil kali $xy(3 - x - y)$ atas
ketiga ``bagian'' $3$ terbesar ketika bagiannya sama
--- yaitu bayangan dua variabel bagi asas AM--GM.
\end{example}

\begin{omfigure}
\begin{tikzpicture}
\begin{axis}[omaxis, width=7.6cm, height=6.6cm,
    xmin=-1.6, xmax=2.1, ymin=-1.6, ymax=2.1,
    xtick={-1,1}, ytick={-1,1},
    xlabel={$x$}, ylabel={$y$}]
  % level curves of f = x^3+y^3-3xy: f = -1 (min), f=-0.5, 0, 1
  \addplot[omDef, thick, domain=-1.35:1.9, samples=80]
    {x^2} node[pos=0.94, left, font=\small]
    {$\frac{\partial f}{\partial x} = 0$};
  \addplot[omThm, thick, domain=-1.35:1.9, samples=80]
    ({x^2}, {x}) node[pos=0.97, below right, font=\small]
    {$\frac{\partial f}{\partial y} = 0$};
  \fill (axis cs:0,0) circle (1.8pt)
    node[below left, font=\small] {pelana};
  \fill (axis cs:1,1) circle (1.8pt)
    node[below right, font=\small] {min};
\end{axis}
\end{tikzpicture}

{\small Kedua kurva $\nabla f = 0$ bagi $f = x^3 + y^3 - 3xy$
berpotongan pada \omterm{thm:b1:multivar:critical}{titik kritis} $(0,0)$ (pelana) dan $(1,1)$
(minimum lokal).}
\end{omfigure}

\begin{remark}[Jebakan yang lazim]
\label{rem:b1:multivar:pitfalls}
\emph{\omterm{def:b1:multivar:partial}{Turunan parsial} boleh ada pada sebuah titik
ketakkontinuan}: karena perangkap radial $g(x,y) = \frac{xy}{x^2+y^2}$ pada
\cref{ex:b1:multivar:trap} mempunyai $\frac{\partial g}{\partial
x}(0,0) = \frac{\partial g}{\partial y}(0,0) = 0$ (karena kedua pembatasan
sumbunya lenyap secara identik), namun $g$ tak berlimit di
titik asalnya --- sebab \omterm{def:b1:multivar:partial}{turunan parsial} hanya menyelidik dua arah, sedangkan \omterm{def:b1:continuity:continuous}{kekontinuan} menuntut
semuanya; dan hanya hipotesis $C^1$ yang memulihkan ketertibannya
(\cref{thm:b1:multivar:c1}).
\emph{Limit garis tak pernah mencukupi}: fungsi pada \cref{exo:b1:multivar:3}
berlimit $0$ sepanjang setiap garis namun tetap tak berlimit ---
jadi selalu cobalah parabola (atau batas kutub yang sah secara seragam pada
$\theta$).
\emph{Kritis itu perlu, bukan cukup}: karena pelana bertaburan
(tiga dari empat titik pada \cref{ex:b1:multivar:fourpoints}); dan
teoremanya berlaku pada \omterm{def:b1:logic:sets}{himpunan} yang \emph{\omterm{def:b1:topology:open}{terbuka}} saja --- karena pada cakram \omterm{def:b1:topology:closed}{tertutup},
ekstremumnya boleh duduk pada batasnya dengan \omterm{def:b1:multivar:partial}{gradien} taknol
(\cref{exo:b1:multivar:12}).
\emph{Kasus bungkam $rt - s^2 = 0$ memang sungguh bungkam}: karena $x^4
+ y^4$ (yang berminimum) dan $x^3 + y^3$ (yang tak keduanya) sama-sama mempunyai $r = s = t
= 0$ di titik asalnya; sehingga hanya kajian tanda yang langsung yang memutuskan
(\cref{exo:b1:multivar:6}, fungsi $h$).
\emph{\omterm{def:b1:multivar:partial}{Gradiennya} \omterm{def:b1:euclid:orthogonal}{ortogonal} terhadap kurva aras, bukan sepanjang
kurvanya}: jadi untuk menyusuri sebuah garis kontur, bergeraklah tegak lurus $\nabla
f$; sedangkan untuk memanjat paling cepat, bergeraklah sepanjangnya --- dan mencampur keduanya membalikkan
geometri setiap peta kontur.
\end{remark}

\begin{remark}[Ke mana dua variabel menuntun]
\label{rem:b1:multivar:whereused}
Bab ini sebuah pintu. \omterm{def:b1:multivar:partial}{Gradien} dan kaidah rantainya meluas
secara harfiah ke $n$ variabel pada jilid Tahun~2, tempat
$o(\norm{(h,k)})$ pada \cref{thm:b1:multivar:c1} menjadi
\emph{diferensialnya} dan uji Monge dibuktikan selengkapnya lewat
rumus Taylor orde dua dan bentuk kuadratik. Adapun kasus khusus
yang dapat diselesaikan tahun \emph{ini} --- yaitu fungsi kuadratik,
yang penjabaran orde duanya persis --- merupakan pokok
soal akhir pekannya, dan kebetulan itulah kasus yang menjalankan
pencocokan data sedunia: yakni regresi \omterm{pb:b1:multivar:1}{kuadrat terkecil}. Sedangkan ekstremum
berkendala (\cref{exo:b1:multivar:5} adalah pratinjaunya) menjadi pengganda
Lagrange pada Tahun~2; dan fungsi harmonik
(\cref{exo:b1:multivar:7}) kembali pada analisis kompleks jilid
Tahun~3.
\end{remark}

\begin{remark}[Cakrawala di dalam Buku 3: tahunnya, ditutup]
\label{rem:b1:multivar:perspectives}
Bab inilah tempat kedua paruh jilid ini berjabat tangan. Adapun
paruh analisisnya memasok perkakasnya satu \omterm{def:b1:derivative:def}{turunan} setiap kali:
karena teorema nilai rata-rata menggerakkan \cref{thm:b1:multivar:c1}, penjabaran
Taylor menggerakkan uji ekstremumnya, dan $\varepsilon$ pada
\cref{ch:b1:topology} kembali dengan bola menggantikan
\omterm{prop:b1:reals:intervals}{selang}. Sedangkan paruh aljabarnya memasok geometrinya: karena \omterm{def:b1:multivar:partial}{gradiennya}
dibaca lewat \omterm{def:b1:euclid:def}{hasil kali dalam} \cref{ch:b1:euclid}
(karena Cauchy--Schwarz menjadikannya arah tercuram), data Monge
$(r, s, t)$ merupakan matriks setangkup
\cref{ch:b1:matrices} dengan uji determinan
\cref{ch:b1:det}, dan soal akhir pekannya menjalankan \omterm{thm:b1:euclid:projection}{proyeksi
ortogonal} pada vektor data. Bahkan kurva
\cref{ch:b1:curves} kembali sebagai \omterm{def:b1:logic:sets}{himpunan} aras. Jadi pembaca yang dapat
membangun ulang mengapa masing-masing kelima serah terima ini bekerja, telah
sesungguhnya mengulang seluruh tahunnya --- dan itulah maksud yang sebenarnya
pada bab \omterm{def:b1:topology:closure}{penutup} ini.
\end{remark}

\begin{omfigure}
\begin{tikzpicture}[scale=1.05]
  \draw[->, thin] (-0.5,0) -- (3.7,0) node[below] {$x$};
  \draw[->, thin] (0,-0.5) -- (0,4.6) node[left] {$y$};
  \draw[omThm, very thick] (-0.28,-0.492) -- (3.3,4.52)
    node[right, font=\small] {$y = 1.4x - 0.1$};
  \fill (0,0) circle (0.06);
  \fill (1,1) circle (0.06);
  \fill (2,3) circle (0.06);
  \fill (3,4) circle (0.06);
  \draw[omDef, dashed, thick] (0,0) -- (0,-0.1);
  \draw[omDef, dashed, thick] (1,1) -- (1,1.3);
  \draw[omDef, dashed, thick] (2,3) -- (2,2.7);
  \draw[omDef, dashed, thick] (3,4) -- (3,4.1);
\end{tikzpicture}

{\small \omterm{pb:b1:multivar:1}{Kuadrat terkecil} dalam satu gambar: yaitu empat titik data,
\omterm{pb:b1:multivar:1}{garis regresinya} $y = 1.4x - 0.1$, dan sisaan tegaknya
(yang putus-putus) yang kuadratnya diminimalkan garisnya --- dengan total $0.2$, yakni
yang terkecil yang dapat dicapai. Adapun soal akhir pekannya menghitung garis ini,
membuktikan bahwa ialah pemberi minimum yang tunggal, dan menyamakan seluruh
konstruksinya dengan sebuah \omterm{thm:b1:euclid:projection}{proyeksi ortogonal} di dalam $\R^4$.}
\end{omfigure}

\section{Latihan}

\begin{exercise}[$\star$]\label{exo:b1:multivar:1}
Hitunglah \omterm{def:b1:multivar:partial}{turunan parsialnya}:
$f(x,y) = x^2 y + \eu^{xy}$; $\;g(x,y) = \ln(x^2 + y^2)$ (pada
$\R^2\setminus\{0\}$); $\;h(x,y) = \arctan\frac yx$ (pada $x > 0$).
\end{exercise}

\begin{exercise}[$\star$]\label{exo:b1:multivar:2}
Pelajarilah \omterm{def:b1:continuity:continuous}{kekontinuannya} di $(0,0)$ (dengan nilai $0$ di situ) bagi:
\[
f(x,y) = \frac{x^2 y}{x^2 + y^2},
\qquad
g(x,y) = \frac{xy}{x^2 + y^2},
\qquad
h(x,y) = \frac{x^3 + y^3}{x^2 + y^2}.
\]
\emph{(\omterm{prop:b1:vspaces:coordinates}{Koordinat} kutub $x = \rho\cos\theta$, $y = \rho\sin\theta$
menolong: batasilah oleh fungsi atas $\rho$ belaka bilamana mungkin.)}
\end{exercise}

\begin{exercise}[$\star\star$]\label{exo:b1:multivar:3}
Misalkan $f(x,y) = \dfrac{x^2 y}{x^4 + y^2}$ ($f(0,0) = 0$). Buktikan bahwa
$f$ berlimit $0$ di titik asalnya \emph{sepanjang setiap garis lurus},
tetapi bahwa $f\bigl(x, x^2\bigr) = \frac12$: sehingga $f$ tak \omterm{def:b1:continuity:continuous}{kontinu} di
$(0,0)$.
\end{exercise}

\begin{exercise}[$\star$]\label{exo:b1:multivar:4}
Periksalah teorema Schwarz dengan tangan pada $f(x, y) = x^3 y^2 + \sin(xy)$.
\end{exercise}

\begin{exercise}[$\star\star$]\label{exo:b1:multivar:5}
Misalkan $f$ berkelas $C^1$ pada $\R^2$ dan $g(t) = f(\cos t, \sin t)$. Nyatakanlah
$g'(t)$ lewat kaidah rantainya. Simpulkanlah bahwa $f$ yang dibatasi ke lingkaran
satuan mencapai ekstremumnya pada titik tempat $\nabla f$ sejajar
vektor jari-jarinya.
\end{exercise}

\begin{exercise}[$\star\star$]\label{exo:b1:multivar:6}
Carilah lalu golongkanlah \omterm{thm:b1:multivar:critical}{titik kritis} pada:
\[
f(x, y) = x^2 + xy + y^2 - 3x,
\qquad
g(x, y) = x^2 - y^2 + 4y,
\qquad
h(x, y) = x^4 + y^4 - 2(x - y)^2 .
\]
\emph{(Bagi $h$, uji determinannya bungkam di titik asalnya:
periksalah $h(x, x)$ dan $h(x, -x)$.)}
\end{exercise}

\begin{exercise}[$\star\star$]\label{exo:b1:multivar:7}
Sebuah fungsi $f$ disebut \emph{harmonik} ketika $\frac{\partial^2 f}{\partial
x^2} + \frac{\partial^2 f}{\partial y^2} = 0$. Periksalah bahwa $x^2 -
y^2$, $xy$, $\eu^x\cos y$ dan $\ln(x^2 + y^2)$ (di luar titik asalnya)
harmonik.
\end{exercise}

\begin{exercise}[$\star\star\star$]\label{exo:b1:multivar:8}
Carilah titik bidang $\{x + 2y - z = 4\} \subseteq \R^3$
yang terdekat dengan titik asalnya, dengan dua cara: lewat \omterm{thm:b1:euclid:projection}{proyeksi ortogonal}
(\cref{ch:b1:euclid}), dan dengan meminimalkan fungsi dua variabel
$f(x, y) = x^2 + y^2 + (x + 2y - 4)^2$ yang diperoleh dengan menghapus $z$.
\end{exercise}

\begin{exercise}[$\star\star\star$]\label{exo:b1:multivar:9}
(Laplacian dalam \omterm{prop:b1:vspaces:coordinates}{koordinat} kutub, sentuhan pertama) Misalkan $f$ berkelas $C^2$ pada
$\R^2 \setminus \{0\}$ dan bersifat \emph{radial}: $f(x, y) =
\varphi\bigl(\sqrt{x^2+y^2}\bigr)$ dengan $\varphi$ berkelas $C^2$ pada
$\intoo{0}{+\infty}$. Buktikan bahwa
\[
\frac{\partial^2 f}{\partial x^2} + \frac{\partial^2 f}{\partial
y^2} = \varphi''(\rho) + \frac{\varphi'(\rho)}{\rho},
\qquad \rho = \sqrt{x^2 + y^2},
\]
lalu carilah semua fungsi harmonik yang radial pada bidang yang tertusuk.
\end{exercise}

\begin{exercise}[$\star\star$]\label{exo:b1:multivar:10}
Berikanlah persamaan bidang singgung grafik $z = xy$ di
titik $(1, 1, 1)$. Lalu carilah semua titik grafik
$f(x,y) = x^3 + y^3 - 3xy$ tempat bidang singgungnya
\emph{mendatar}, lalu kaitkanlah jawabannya dengan
\cref{ex:b1:multivar:extrema}.
\end{exercise}

\begin{exercise}[$\star\star$]\label{exo:b1:multivar:11}
Misalkan $f$ berkelas $C^1$ pada $\R^2$.
\begin{enumerate}
  \item Jika $\frac{\partial f}{\partial x} = 0$ di mana-mana, buktikanlah
        bahwa $f(x, y)$ hanya bergantung pada $y$.
  \item Carilah semua penyelesaian $C^1$ atas persamaan
        $\frac{\partial f}{\partial x} = \frac{\partial
        f}{\partial y}$ pada $\R^2$. \emph{(Tetapkanlah $g(u, v) =
        f\bigl(\frac{u+v}2, \frac{u-v}2\bigr)$ lalu hitunglah
        $\frac{\partial g}{\partial v}$ lewat kaidah rantainya.)}
\end{enumerate}
\end{exercise}

\begin{exercise}[$\star\star\star$]\label{exo:b1:multivar:12}
Carilah maksimum dan minimum global $f(x, y) = xy$ pada
cakram \omterm{def:b1:topology:closed}{tertutup} $x^2 + y^2 \leq 1$. \emph{(Akuilah teorema nilai
ekstrem dua variabelnya: bahwa fungsi \omterm{def:b1:continuity:continuous}{kontinu} pada cakram \omterm{def:b1:topology:closed}{tertutup}
mencapai batasnya --- yang dibuktikan pada jilid Tahun~2. Perlakukanlah
cakram \omterm{def:b1:topology:open}{terbukanya} lewat \omterm{thm:b1:multivar:critical}{titik kritis} dan lingkaran batasnya lewat
pemarameteran \cref{exo:b1:multivar:5}.)}
\end{exercise}

\section{Soal: kuadrat terkecil dan garis regresi}

\begin{problem}\label{pb:b1:multivar:1}
Diberikan $n$ titik data, garis lurus yang mana yang lewat ``paling dekat''
dengan semuanya? Legendre dan Gauss menjawab: yaitu garis yang meminimalkan
jumlah galat tegak yang \emph{dikuadratkan} --- karena peminimalan
itulah yang persis terpecahkan oleh aljabar \omterm{def:b1:linmaps:def}{linear}. Soal ini
mula-mula membuktikan uji Monge \cref{met:b1:multivar:monge}
secara jujur bagi fungsi kuadratik (yakni satu-satunya kasus tempat
penjabaran orde duanya persis), lalu membangun \omterm{pb:b1:multivar:1}{persamaan
normal}, \omterm{pb:b1:multivar:1}{garis regresi}, dan koefisien korelasinya
di atas geometri Euklides
\cref{ch:b1:euclid}.\index{kuadrat terkecil}\index{garis regresi}
\index{persamaan normal}\index{korelasi}
\medskip
\noindent\textbf{Bagian I --- Fungsi kuadratik: uji Monge,
dibuktikan.} Tetapkanlah bilangan real $r, s, t$ lalu misalkan $q(h, k) = r h^2 + 2s hk +
t k^2$.
\begin{enumerate}
  \item Andaikan $r \neq 0$. Tegakkanlah bentuk kuadrat lengkapnya
        \[
        q(h, k) = r\Bigl(h + \frac{s}{r}k\Bigr)^{\!2}
        + \frac{rt - s^2}{r}\,k^2 ,
        \]
        lalu simpulkanlah: bahwa jika $rt - s^2 > 0$, maka $q$ bertanda tegas
        seperti $r$ di luar titik asalnya; sedangkan jika $rt - s^2 < 0$, maka $q$ mengambil kedua
        tandanya.
  \item Selesaikanlah kasus sisanya: $r = 0$, $t \neq 0$
        (yang setangkup); dan $r = t = 0$ ($q = 2shk$). Simpulkanlah:
        bahwa $q$ mengambil kedua tandanya jika dan hanya jika $rt - s^2 < 0$, dan bahwa $q$
        lenyap hanya di titik asalnya jika dan hanya jika $rt - s^2 > 0$.
  \item Kini misalkan $f(x, y) = \frac12\bigl(r x^2 + 2s xy + t
        y^2\bigr) + \beta x + \gamma y + \delta$ sebuah fungsi
        kuadratik yang bertitik kritis $X_0 = (a, b)$. Buktikan
        penjabaran yang \emph{persis}
        \[
        f(X_0 + (h,k)) = f(X_0) + \tfrac12\,q(h, k)
        \qquad (\text{tanpa sisa}),
        \]
        lalu simpulkanlah penggolongan Monge bagi fungsi
        kuadratik: yaitu minimum global yang tegas jika $rt - s^2 > 0$, $r >
        0$; maksimum global yang tegas jika $rt - s^2 > 0$, $r < 0$;
        dan pelana jika $rt - s^2 < 0$.
  \item Andaikan $rt - s^2 > 0$ dan $r > 0$ (sehingga juga $t > 0$).
        Buktikan batas keko\-ersifan yang eksplisit
        \[
        q(h, k) \geq c\,(h^2 + k^2),
        \qquad
        c = \frac{rt - s^2}{r + t} > 0 .
        \]
        \emph{(Tunjukkanlah bahwa bentuk kuadratik $q - c(h^2 + k^2)$,
        yang berkoefisien $r - c$, $s$, $t - c$, masih mempunyai
        data uji diskriminan yang taknegatif.)}
  \item Simpulkanlah bahwa fungsi kuadratik yang bagian kuadratiknya tegas
        positif menuju $+\infty$ ketika $\norm{(x,y)} \to
        \infty$, dan karenanya mempunyai pemberi minimum \emph{global
        yang tunggal}: yaitu titik kritisnya. Bandingkanlah dengan
        \cref{ex:b1:multivar:extrema}, tempat minimum lokal sebuah
        fungsi yang takkuadratik ternyata tak global.
\end{enumerate}

\medskip
\noindent\textbf{Bagian II --- \omterm{pb:b1:multivar:1}{Persamaan normalnya}.} Misalkan $C_1,
C_2, b \in \R^n$ (dengan \omterm{def:b1:euclid:def}{hasil kali dalam} kanoniknya), dan
\[
f(u, v) = \norm{u\,C_1 + v\,C_2 - b}^2 .
\]
\begin{enumerate}[resume]
  \item Jabarkanlah $f$ lalu tunjukkanlah bahwa ia fungsi kuadratik atas $(u,
        v)$ yang koefisien bagian kuadratiknya $r = 2\norm{C_1}^2$,
        $s = 2\langle C_1, C_2\rangle$, $t = 2\norm{C_2}^2$.
  \item Tunjukkan bahwa $rt - s^2 > 0$ jika dan hanya jika $(C_1, C_2)$
        \omterm{def:b1:vspaces:free}{bebas} \emph{(yaitu kasus kesamaan Cauchy--Schwarz,
        \cref{thm:b1:euclid:cs})}; andaikanlah ini mulai sekarang.
  \item Tunjukkan bahwa persamaan kritisnya $\nabla f = 0$ adalah
        \emph{\omterm{pb:b1:multivar:1}{persamaan normalnya}}
        \[
        \begin{pmatrix}
        \norm{C_1}^2 & \langle C_1, C_2\rangle\\
        \langle C_1, C_2\rangle & \norm{C_2}^2
        \end{pmatrix}
        \begin{pmatrix} u\\ v\end{pmatrix}
        =
        \begin{pmatrix} \langle C_1, b\rangle\\ \langle C_2,
        b\rangle\end{pmatrix}
        \]
        --- dengan matriks Gram \cref{exo:b1:euclid:11} di
        kiri --- dan bahwa ia berkata tepat: $b - (uC_1 + vC_2)
        \perp C_1, C_2$. Simpulkanlah bersama Bagian I dan
        \cref{thm:b1:euclid:projection}: bahwa pemberi minimum yang tunggal
        memberikan $p = uC_1 + vC_2 =$ \omterm{thm:b1:euclid:projection}{proyeksi ortogonal}
        $b$ ke $\operatorname{Vect}(C_1, C_2)$.
  \item Tunjukkan bahwa nilai minimalnya adalah $\norm{b}^2 - \langle p,
        b\rangle = \norm{b - p}^2$, lalu gambarlah
        gambaran Pythagorasnya.
\end{enumerate}

\medskip
\noindent\textbf{Bagian III --- \omterm{pb:b1:multivar:1}{Garis regresinya}.} Dengan titik data
$(x_1, y_1), \dots, (x_n, y_n)$, yang $x_i$-nya tak semua sama. Minimalkanlah
\[
E(\alpha, \beta) = \sum_{i=1}^{n}\bigl(y_i - \alpha - \beta
x_i\bigr)^2 .
\]
Tulislah $\bar x = \frac1n\sum x_i$, $\bar y = \frac1n\sum y_i$,
$v_x = \frac1n\sum x_i^2 - \bar x^2$, $v_y = \frac1n\sum y_i^2 -
\bar y^2$, $c_{xy} = \frac1n\sum x_i y_i - \bar x\bar y$.
\begin{enumerate}[resume]
  \item Kenalilah Bagian II dengan $C_1 = (1, \dots, 1)$, $C_2 =
        (x_1, \dots, x_n)$, $b = (y_1, \dots, y_n)$, lalu tulislah
        \omterm{pb:b1:multivar:1}{persamaan normalnya}
        \[
        n\alpha + \Bigl(\sum x_i\Bigr)\beta = \sum y_i,
        \qquad
        \Bigl(\sum x_i\Bigr)\alpha + \Bigl(\sum
        x_i^2\Bigr)\beta = \sum x_i y_i .
        \]
  \item Pecahkanlah:
        \[
        \beta = \frac{c_{xy}}{v_x},
        \qquad
        \alpha = \bar y - \beta\,\bar x ,
        \]
        lalu amatilah bahwa \emph{\omterm{pb:b1:multivar:1}{garis regresinya}} $y = \alpha
        + \beta x$ lewat titik reratanya $(\bar x, \bar
        y)$.
  \item Periksalah $v_x > 0$ tepat ketika $x_i$ tak semuanya
        sama, lalu cocokkanlah ini dengan pertanyaan 7.
  \item Buktikan bahwa galat minimalnya adalah
        \[
        E_{\min} = n\Bigl(v_y - \frac{c_{xy}^2}{v_x}\Bigr)
        = n\,v_y\,(1 - \rho^2),
        \qquad
        \rho = \frac{c_{xy}}{\sqrt{v_x v_y}}
        \quad (v_y \neq 0).
        \]
        Simpulkanlah $\abs\rho \leq 1$, dengan $\abs\rho = 1$ jika dan
        hanya jika datanya sejajar sempurna.
  \item Contoh yang dikerjakan: bagi data $(0,0)$, $(1,1)$, $(2,3)$,
        $(3,4)$, hitunglah $\bar x, \bar y, v_x, c_{xy}$,
        \omterm{pb:b1:multivar:1}{garis regresinya}, keempat sisaannya dan $E_{\min}$.
  \item Buktikan bahwa sisaan $\varepsilon_i = y_i - \alpha -
        \beta x_i$ pada garis yang optimal memenuhi $\sum_i
        \varepsilon_i = 0$ dan $\sum_i x_i\varepsilon_i = 0$,
        lalu tafsirkanlah keduanya lewat keortogonalannya.
\end{enumerate}

\medskip
\noindent\textbf{Bagian IV --- Ragam dan penerapannya.}
\begin{enumerate}[resume]
  \item (Tetapan terbaik) Tunjukkan bahwa tetapan $\alpha$ yang
        meminimalkan $\sum_i (y_i - \alpha)^2$ adalah reratanya $\bar
        y$, dan bahwa nilai minimalnya $n\,v_y$: jadi variansnya
        mengukur kegagalan datanya menjadi tetap.
  \item (Garis lewat titik asal) Tunjukkan bahwa kemiringan yang meminimalkan
        $\sum_i (y_i - \beta x_i)^2$ adalah $\beta_0 = \frac{\sum
        x_iy_i}{\sum x_i^2}$, lalu berikanlah syarat pada datanya
        agar $\beta_0$ berimpit dengan kemiringan
        $c_{xy}/v_x$ pada pertanyaan 11.
  \item (Jarak ke sebuah garis, lagi) Bagi sebuah titik $M$ dan garis
        $D$ lewat $P_0$ yang berarah vektor satuan $w$,
        minimalkanlah $g(\tau) = \norm{P_0 + \tau w - M}^2$ lalu
        simpulkanlah $d(M, D)^2 = \norm{M - P_0}^2 - \langle M - P_0,
        w\rangle^2$; lalu pulihkanlah rumus $d = \frac{\abs{a x_0 +
        b y_0 + c}}{\sqrt{a^2 + b^2}}$ bagi garis $ax + by + c
        = 0$ pada bidangnya.
  \item (Analisis varians) Buktikan penguraian $v_y =
        \beta^2 v_x + \frac{E_{\min}}n$: bahwa varians
        $y_i$ terbelah menjadi bagian yang dijelaskan garisnya ditambah
        varians sisaannya.
  \item (Pencocokan parabolik) Untuk mencocokkan $y = a + bx + cx^2$, tunjukkanlah bahwa
        \omterm{pb:b1:multivar:1}{persamaan normalnya} adalah sistem $3 \times 3$ dengan
        matriks momen
        \[
        \begin{pmatrix}
        n & \sum x_i & \sum x_i^2\\
        \sum x_i & \sum x_i^2 & \sum x_i^3\\
        \sum x_i^2 & \sum x_i^3 & \sum x_i^4
        \end{pmatrix},
        \]
        yaitu matriks Gram yang dapat dibalik segera setelah tiga di antara
        $x_i$-nya berbeda \emph{(lewat kebebasan $(1, X, X^2)$
        yang dicuplik pada datanya, \cref{exo:b1:euclid:11}; bandingkanlah
        matriks momen pada soal akhir pekan
        \cref{ch:b1:det})}.
\end{enumerate}

\medskip
\noindent\textbf{Bagian V --- Ketegaran, dan sintesisnya.}
\begin{enumerate}[resume]
  \item Golongkanlah \omterm{thm:b1:multivar:critical}{titik kritis} ketiga fungsi kuadratik
        berikut
        \[
        x^2 + xy + y^2, \qquad x^2 + 3xy + y^2, \qquad
        x^2 + 2xy + y^2 ,
        \]
        lewat Bagian I, dengan memperlakukan kasus ketiganya yang merosot dengan tangan
        (di manakah minimumnya tercapai?).
  \item (Percobaan pencilan) Tambahkanlah titik $(10, 0)$ pada
        data pertanyaan 14 lalu hitunglah ulang kemiringannya $\beta$.
        Apakah yang terjadi, dan mengapakah galat yang \emph{dikuadratkan} begitu
        peka terhadap satu titik yang jauh?
  \item (\omterm{pb:b1:multivar:1}{Kuadrat terkecil} terbobot) Diberikan bobot $w_i > 0$,
        minimalkanlah $\sum_i w_i(y_i - \alpha - \beta x_i)^2$: tunjukkanlah
        bahwa penyelesaiannya diberikan rumus yang sama seperti pertanyaan
        11 dengan rerata, varians dan kovarians yang terbobot
        (definisikanlah semuanya).
  \item Tunjukkan bahwa bagi garis yang optimal, $E_{\min} = 0$ memaksa
        semua titiknya terletak persis padanya, lalu kaitkanlah dengan
        kasus kesamaan $\abs\rho = 1$ pada pertanyaan 13; lalu ujilah pada
        data yang sejajar $(1,1), (2,3), (3,5)$.
  \item Sintesis, dalam empat kalimat: mengapakah fungsi kuadratik
        merupakan satu-satunya kelas tempat uji orde dua bab ini
        tak memerlukan teorema yang diakui; bagaimanakah \omterm{pb:b1:multivar:1}{persamaan normalnya}
        menyamakan peminimalan analisisnya dengan \omterm{thm:b1:euclid:projection}{proyeksi
        ortogonal} \cref{ch:b1:euclid}; apakah yang diukur koefisien
        korelasinya dan ketaksamaan yang mana yang membatasinya;
        dan bagian yang mana pada Bab 18--23 (yaitu \omterm{def:b1:vspaces:free}{basis}, matriks Gram dan
        matriks momen, \omterm{def:b1:linmaps:projection}{proyeksi}) yang muncul kembali. Namailah
        metode dan persamaan yang dipelajari pada Bagian II--III.

\end{enumerate}
\end{problem}
